\labeledsection{Introduction}
The topic of this research report is rather utilizing and implementing Dependency-Injection in Java-Applications comes with more potential dangers than benefits or not.

Overall, this report will first explain what Dependency-Injection in Java looks lie, what is important to know and how it works on a basic level.

Overall, this report will first explain what Dependency-Injection in Java looks like, what is needed for it and how it works.
Then, a qualitative multi-stage experiment is discussed and finally the results and findings are presented.

\begin{enumerate}
    \item What is Java Dependency-Injection?
    \item What kind of Java Dependency-Injection are there?
      'wanted dependency injection' (Plugin-System)
      'abuse of wanted Dependency-Injection'
      'unwanted Dependency-Injection' [-> only mention]
    \item What are the positive benefits vs. negative impacts?c
\end{enumerate}

\chapter{The experiment}
    % Research topic
    The experiment will prove the potential risk of actively using Java-Dependency-Injection to add or modify data, structures and objects at runtime.
    % Research element
    During the experiment a made-up application will be exposed to three beneficial exploits.
    With each iteration the security of said application will improve and shield against the previous exploit.  
    This way, different security levels and potential risks are measured. 
    Additionally, we expect the total effective damage to decrease with each security level.
    % Variables
    \begin{tabular}{|r|l|}
        /hline
        Security level & The implemented security measures, placed on an artificial scale. Generally, with each incrementum the security of the overall application is improved significantly. \\\hline
        Effective damage & Damage to the application or business process involved in using this application, which is \textbf{measurable} and can be presented \\\hline
        Potential risk & Risks that are present and could be exploited by others, though \textbf{not} directly measurable and only hypothetical. \\\hline
    \end{tabular}
    % Values ... Measure!